\section{Problem formulation}\label{sec:problem}
\subsection{Replanning snapshot}
A central planner is called whenever a replanning trigger fires (Section~\ref{sec:triggers}). Each UAV keeps the task
that it is currently flying to or serving. This committed task determines when and where the UAV becomes available
again and how much energy it has left. Each open task has a location, a release time, a deadline by which its service
must be completed, a service duration and a reward. A plan assigns an ordered sequence of open tasks to each UAV, and
tasks may remain unassigned. The planning problem is therefore a prize-collecting routing problem with time
windows~\cite{solomon1987algorithms,vansteenwegen2011orienteering}. Arrival, start and finish times along a route are
computed in the usual way for time windows~\cite{solomon1987algorithms,savelsbergh1992vrptw}. A route is feasible if
every service ends before its deadline and the UAV returns to the depot with an energy reserve of 10\% of its capacity.
Flight consumes energy in proportion to the distance, hovering during service consumes energy over time, and waiting
consumes no energy. A plan $\pi$ is evaluated by
\begin{equation}
J(\pi)=\alpha R(\pi)-\gamma E(\pi)-\delta C(\pi),
\end{equation}
where $R$ is the reward of the tasks served on time, $E$ is the flown distance in kilometres, and $C$ is the number of
tasks that move to a different UAV than in the previous plan. The weights are $\alpha=1$, $\gamma=0.05$ per kilometre
and $\delta=0.3$. Plan quality is reported in percent of $W$, the total reward of the open tasks.

\subsection{Publication delay}
Computation is not free. A plan that is computed for $\tau$ milliseconds is released $\kappa\tau$ milliseconds later in
mission time. The latency factor $\kappa$ converts planning time into mission time. The value $\kappa=1$ corresponds to
a planner that runs in real time on the evaluation computer, whereas larger values represent slower on-board
processors. A UAV that is already idle must wait for the new plan, so its whole route starts later, and tasks whose
deadline can no longer be met are lost. With $J_{\mathrm{eff}}(\pi,\tau)$ as the objective of $\pi$ after this shift,
the net utility of publishing $\pi$ after $\tau$ milliseconds is
\begin{equation}
U(\pi,\tau)=J_{\mathrm{eff}}(\pi,\tau)-\mu W\tau,
\end{equation}
where $\mu$ is the price of one millisecond of computation as a fraction of $W$. The planner thus chooses both a plan
and the moment at which it stops computing.

\subsection{Uncertain execution and dissemination}\label{sec:uncertainty}
The planner uses a nominal model: straight legs at the cruise speed of each UAV and nominal service durations. The
executed mission differs from this model in two ways. First, the finish time of task $j$ is $F_j=f_j+e_j$, where $f_j$
is the nominal finish time computed at publication and $e_j$ is an execution error caused by acceleration, wind,
speed tracking and service-time variability. Second, the plan reaches the vehicles after an additional dissemination
latency $L$ that is not under the control of the planner. A task yields its reward only if its actual finish time meets
its deadline, which may have been tightened in the meantime. Following common practice in dispatching, a UAV skips a
task at departure if the nominal estimate from its actual state already misses the deadline, and the task returns to
the pool of open tasks. The performance measure of a mission is its utility, the reward of all tasks served on time
divided by the total reward of all tasks of the mission.

\subsection{Replanning triggers}\label{sec:triggers}
Three trigger policies are considered. The event trigger replans on task arrivals, UAV failures and deadline changes
only. The periodic trigger additionally replans every $T$ seconds of mission time. The deviation trigger additionally
replans when the actual finish of a planned task deviates from its planned finish by more than $\theta$ seconds. During
planning, every UAV finishes its committed task and then waits for the new plan, and events that arrive during planning
are queued. These semantics are identical in both simulation tiers.
